\chapter{Interview Cheat Sheets}
\label{app:interview-cheatsheets}

Printable one-pagers for fast review before an interview. Each page is a memory
trigger, not a substitute for the chapter. Read the goal, walk the checklist,
say the answer out loud, then use the pressure pivot.

\small

\section{Book overview}
\label{sec:cheat-overview}

\textit{Universal rule.} State the decision, name the evidence, explain the
tradeoff, and say what you would do next.

\paragraph{Lifecycle spine.}
Model the link $\rightarrow$ choose the architecture $\rightarrow$ control
wavelength and margins $\rightarrow$ establish readiness $\rightarrow$ qualify
life $\rightarrow$ validate manufacturing $\rightarrow$ operate the fleet
$\rightarrow$ investigate failures.

\begin{center}
\footnotesize
\begin{tabular}{@{}cp{0.72\linewidth}@{}}
\toprule
Ch & Main question \\
\midrule
4 & What sets BER and margin? \\
5 & How do the source and modulator behave? \\
6 & How do wavelength and WDM stay controlled? \\
7 & What evidence is needed before release? \\
8 & Will time and exposure cause permanent degradation? \\
9 & Can the factory build and control it repeatedly? \\
10 & How does optical behavior affect the system and workload? \\
11 & How do I find and prevent the real cause of a failure? \\
\bottomrule
\end{tabular}
\end{center}

\paragraph{How to use these sheets.}
Spend about 2--3 minutes per page. Chapters~7--11 first if time is short. Drill
the rapid-fire flashcards at the end of this appendix in one-breath answers.
Full depth:
\Cref{ch:models,ch:lasers,ch:wdm,ch:product-readiness,ch:reliability,ch:manufacturing,ch:networking,ch:failure-modes}.

\clearpage
\section{Chapter 4: Quantitative link behavior}
\label{sec:cheat-ch4}

\textit{Goal.} Explain what sets bit-error rate (BER) and margin without hiding
behind a single number.

\paragraph{Mental checklist.}
Define planes $\rightarrow$ Build budget $\rightarrow$ Add noise $\rightarrow$
Find boundary $\rightarrow$ Measure margin $\rightarrow$ Check assumptions

\paragraph{What I need to know.}
\begin{itemize}
\item BER is an outcome of signal, noise, bandwidth, distortion, and decision
      behavior.
\item Average power alone does not define receiver performance.
\item A waterfall shows how performance changes as margin is removed.
\item A floor suggests that adding more power may not solve the problem.
\item Always state the reference plane and forward-error correction (FEC)
      condition.
\end{itemize}

\paragraph{Key distinctions.}
\begin{center}
\footnotesize
\begin{tabular}{@{}ll@{}}
\toprule
Do not confuse & With \\
\midrule
Average optical power & Receiver sensitivity at a named BER \\
Waterfall shift & Error floor \\
Pre-FEC activity & Post-FEC residual \\
\bottomrule
\end{tabular}
\end{center}

\paragraph{Numbers worth remembering.}
Margin is the distance from the operating point to the measured failure
boundary.

\paragraph{Trap.}
Trap: Quoting BER without plane, pattern, or FEC condition.\\
Better: Name the plane and the BER or FEC objective first.

\paragraph{60-second answer.}
``I start by naming the reference plane and the BER or FEC condition. Then I
build a simple budget for signal and the main noise or impairment terms, find
the failure boundary, and measure how much margin remains. If the waterfall
shifts, more power may help. If I see a floor, I look for noise, reflection, or
decision limits instead of chasing average power.''

\paragraph{Pressure.}
Pressure: ``Our average Rx power looks fine. Why is BER bad?''\\
Pivot: ``Stable average power rules out gross loss, not timing, noise, or
distortion. I would pull a waterfall and check for a floor.''

\footnotesize Depth: \Cref{ch:models}.

\clearpage
\section{Chapter 5: Sources and modulation}
\label{sec:cheat-ch5}

\textit{Goal.} Choose and operate a source so the modulated link closes the
system budget, not just a lab power target.

\paragraph{Mental checklist.}
Choose source $\rightarrow$ Set operating point $\rightarrow$ Modulate
$\rightarrow$ Control temperature $\rightarrow$ Check quality $\rightarrow$
Reserve headroom

\paragraph{What I need to know.}
\begin{itemize}
\item Laser output power is not the same as useful modulated signal.
\item Bias, temperature, slope efficiency, extinction, chirp, and bandwidth
      interact.
\item The control loop needs room to respond over temperature and aging.
\item The best source closes system, reliability, and manufacturing budgets,
      not only the best nominal lab result.
\end{itemize}

\paragraph{Key distinctions.}
\begin{center}
\footnotesize
\begin{tabular}{@{}ll@{}}
\toprule
Do not confuse & With \\
\midrule
CW launch power & Modulated OMA / usable signal \\
Nominal lab result & Worst-corner closed budget \\
Bias setting & Remaining control headroom \\
\bottomrule
\end{tabular}
\end{center}

\paragraph{Trap.}
Trap: Picking the laser with the highest room-temperature power.\\
Better: Ask which source still closes margin, reliability, and factory control
at the claimed corners.

\paragraph{60-second answer.}
``I choose the source from the system budget: reach, modulation, temperature,
reliability, and manufacturability. I set bias and temperature with headroom for
aging, then check modulated quality, not only CW power. If the loop is already
at the stop at one corner, the design is not ready.''

\paragraph{Pressure.}
Pressure: ``Why not just raise bias for more power?''\\
Pivot: ``Bias also moves extinction, chirp, reliability, and headroom. I would
show the net margin, not one power number.''

\footnotesize Depth: \Cref{ch:lasers}.

\clearpage
\section{Chapter 6: WDM and wavelength control}
\label{sec:cheat-ch6}

\textit{Goal.} Keep channels on the grid across temperature, process, and aging
with enough control authority left.

\paragraph{Mental checklist.}
Set grid $\rightarrow$ Model drift $\rightarrow$ Sense error $\rightarrow$
Correct wavelength $\rightarrow$ Check authority $\rightarrow$ Handle failure

\paragraph{What I need to know.}
\begin{itemize}
\item Wavelength is a control problem, not only a nominal setting.
\item Temperature, process, aging, and neighboring channels move the operating
      point.
\item Heater or TEC authority must remain at the worst corner.
\item Passing at room temperature does not prove control across life.
\item Monitor both wavelength and how hard the loop is working.
\end{itemize}

\paragraph{Key distinctions.}
\begin{center}
\footnotesize
\begin{tabular}{@{}ll@{}}
\toprule
Do not confuse & With \\
\midrule
Locked flag & Remaining control headroom \\
Room-temp pass & Life and corner control \\
Source drift & Filter or ring drift \\
\bottomrule
\end{tabular}
\end{center}

\paragraph{Trap.}
Trap: Treating ``locked'' as proof of margin.\\
Better: Check lock error and actuator demand at the worst claimed corner.

\paragraph{60-second answer.}
``I treat wavelength as a closed-loop problem. I name the grid, the drift
sources, and the sense/correct path. Then I ask whether the actuator still has
room at the hot and aged corners. A lock bit without headroom is not a passing
control design.''

\paragraph{Pressure.}
Pressure: ``It locks on the bench. Ship it?''\\
Pivot: ``Not until I see authority and wavelength error across temperature and
aging, not only a room-temp lock.''

\footnotesize Depth: \Cref{ch:wdm}.

\clearpage
\section{Chapter 7: Product readiness}
\label{sec:cheat-ch7}

\textit{Goal.} Know what evidence is required before release, and which question
each stage answers.

\paragraph{Mental checklist.}
Define $\rightarrow$ Review $\rightarrow$ Bring up $\rightarrow$ Characterize
$\rightarrow$ Verify/validate $\rightarrow$ Qualify $\rightarrow$ Validate
factory $\rightarrow$ Pilot $\rightarrow$ Ramp $\rightarrow$ Monitor
$\rightarrow$ Learn

\paragraph{What I need to know.}
\begin{itemize}
\item Product readiness is the umbrella lifecycle.
\item Characterization maps behavior; verification checks frozen requirements.
\item System validation proves intended use.
\item Reliability qualification and manufacturing validation answer different
      questions.
\item Later evidence cannot replace missing earlier evidence.
\item Program labels such as EVT, DVT, and PVT do not define what the evidence
      proves.
\end{itemize}

\paragraph{Key distinctions.}
\begin{center}
\footnotesize
\begin{tabular}{@{}ll@{}}
\toprule
Do not confuse & With \\
\midrule
Characterization & Verification \\
System validation & Reliability qualification \\
Manufacturing validation & Qualification life claim \\
Phase label (EVT/DVT/PVT) & Evidence content \\
\bottomrule
\end{tabular}
\end{center}

\paragraph{Trap.}
Trap: Using a phase name as proof of readiness.\\
Better: Translate the phase into hardware, evidence, open risks, and the exit
decision.

\paragraph{60-second answer.}
``I start from the claim and the decision gate. Characterization maps the
behavior. Verification checks frozen requirements. System validation proves
intended use. Qualification and manufacturing validation answer life and factory
questions separately. I will not use a later pass to paper over a missing earlier
step.''

\paragraph{Pressure.}
Pressure: ``We are in PVT. Are we done?''\\
Pivot: ``PVT is a program label. I ask what evidence exists for the remaining
risks and what decision it supports.''

\footnotesize Depth: \Cref{ch:product-readiness,tab:ladder}.

\clearpage
\section{Chapter 8: Reliability qualification}
\label{sec:cheat-ch8}

\textit{Goal.} Build a life or environment claim from mechanism and evidence,
not from a copied test list.

\paragraph{Mental checklist.}
Claim $\rightarrow$ Mechanism $\rightarrow$ Stress $\rightarrow$ Observable
$\rightarrow$ Acceptance $\rightarrow$ Samples $\rightarrow$ Confidence
$\rightarrow$ Decision $\rightarrow$ Handoff

\paragraph{What I need to know.}
\begin{itemize}
\item Start with the claim and failure mechanism, not a test list.
\item The stress must accelerate the field mechanism without creating a
      different one.
\item Separate reversible hot behavior from permanent degradation.
\item Zero failures do not mean zero failure rate.
\item HTOL is qualification evidence; burn-in is an optional production screen.
\item Standards provide methods, not the complete argument.
\end{itemize}

\paragraph{Key distinctions.}
\begin{center}
\footnotesize
\begin{tabular}{@{}ll@{}}
\toprule
Do not confuse & With \\
\midrule
HTOL / qual stress & Production burn-in screen \\
Reversible thermal behavior & Permanent wear-out \\
Component FIT & System availability \\
\bottomrule
\end{tabular}
\end{center}

\paragraph{Numbers worth remembering.}
One FIT $=1$ failure per $10^9$ device-hours. Zero observed fails still leaves a
one-sided upper bound on rate.

\paragraph{Trap.}
Trap: Saying ``we passed GR-468'' as the whole reliability argument.\\
Better: State the claim, mechanism, stress relevance, sample basis, and residual
risk.

\paragraph{60-second answer.}
``I start from the life claim and the mechanism. I choose a stress that
accelerates that mechanism, define the observable and acceptance rule, then size
samples for the confidence I need. Zero fails still leave an upper bound. I
handoff production screens only when the screen separates real escapes.''

\paragraph{Pressure.}
Pressure: ``Can we skip HTOL if burn-in looks clean?''\\
Pivot: ``Burn-in is a screen. It does not replace the qualification claim unless
the argument explicitly says so.''

\footnotesize Depth: \Cref{ch:reliability}.

\clearpage
\section{Chapter 9: Manufacturing validation}
\label{sec:cheat-ch9}
\enlargethispage{3\baselineskip}

\textit{Goal.} Prove the factory can repeatedly build, measure, and control the
released design.

\paragraph{Mental checklist.}
Freeze $\rightarrow$ Represent $\rightarrow$ Trace $\rightarrow$ Trust
measurement $\rightarrow$ Read yield $\rightarrow$ Check stability $\rightarrow$
Measure capability $\rightarrow$ Control $\rightarrow$ FAIR $\rightarrow$ React
$\rightarrow$ Ramp

\paragraph{What I need to know.}
\begin{itemize}
\item Freeze the production reference; changes need approval.
\item Validate the measurement system before believing yield.
\item First-pass yield shows process health; final yield may hide rework.
\item Look at the distribution before using capability numbers.
\item Stable process first, then $C_p$ and $C_{pk}$.
\item FAIR is the controlled first-article evidence package that gates open
      volume after a relevant tooling, site, material, silicon, assembly, test,
      or firmware change; evidence depth follows the affected risk.
\item Every control needs a trigger, owner, action, and restart rule.
\item Passing good units does not prove a screen catches bad ones.
\end{itemize}

\paragraph{Key distinctions.}
\begin{center}
\footnotesize
\begin{tabular}{@{}ll@{}}
\toprule
Do not confuse & With \\
\midrule
Control limits & Specification limits \\
ATP / acceptance limits & Specification limits \\
First-pass yield & Final yield \\
$C_p$ (spread) & $C_{pk}$ (spread and offset) \\
\bottomrule
\end{tabular}
\end{center}

\paragraph{Numbers worth remembering.}
$C_p$ asks whether the process spread could fit the specification if centered.
$C_{pk}$ includes the actual process offset.

\paragraph{Context.}
\begin{center}
\footnotesize
\begin{tabular}{@{}ll@{}}
\toprule
Context & First priority \\
\midrule
Lab / learning build & Learn and separate variables \\
Pilot lot & Protect lot integrity; collect representative evidence \\
Production & Protect shipments; execute the reaction plan \\
\bottomrule
\end{tabular}
\end{center}
Same flow, different first priority.

\paragraph{Trap.}
Trap: Using final yield to hide heavy rework.\\
Better: Separate first-pass yield, retest, rework, scrap, and final yield.

\paragraph{60-second answer.}
``I freeze the production reference, build a representative population, and clear
the measurement system before I trust yield. I look at first-pass yield and the
distribution, establish stability, then interpret $C_p$ or $C_{pk}$. Controls
need owners and reaction plans. I ramp only when evidence supports the next
volume.''

\paragraph{Pressure.}
Pressure: ``Final yield is 99\%. Ship?''\\
Pivot: ``Not until I know first-pass yield and what the recovery path is costing
and hiding.''

\footnotesize Depth:
\Cref{ch:manufacturing,tab:mfg-three-limits,sec:mfg-fair-checklist}.

\clearpage
\section{Chapter 10: Optical links in operation}
\label{sec:cheat-ch10}

\textit{Goal.} Connect physical margin to FEC, recovery, workload impact, and
fleet decisions.

\paragraph{Mental checklist.}
Observe $\rightarrow$ Read FEC $\rightarrow$ Find bursts $\rightarrow$ Check
recovery $\rightarrow$ Use telemetry $\rightarrow$ Compare cohorts $\rightarrow$
Judge impact $\rightarrow$ Contain

\paragraph{What I need to know.}
\begin{itemize}
\item Link-up is not the same as link health.
\item Corrected errors can reveal shrinking margin before traffic fails.
\item Average BER can hide short bursts.
\item Recovery time and repeated retrains matter.
\item Telemetry is useful when it separates a cause or changes an action.
\item Component failure rate is not system availability.
\item At fleet scale, a rare per-link event may happen frequently.
\end{itemize}

\paragraph{Key distinctions.}
\begin{center}
\footnotesize
\begin{tabular}{@{}ll@{}}
\toprule
Do not confuse & With \\
\midrule
Link-up & Link health over time \\
Corrected-error rise & Service already failing \\
Component FIT & Fleet event rate and availability \\
\bottomrule
\end{tabular}
\end{center}

\paragraph{Context.}
\begin{center}
\footnotesize
\begin{tabular}{@{}ll@{}}
\toprule
Context & First priority \\
\midrule
Lab & Learn and reproduce \\
Pilot & Bound exposure; keep observability \\
Production fleet & Protect the workload; contain first \\
\bottomrule
\end{tabular}
\end{center}
Same flow, different first priority.

\paragraph{Trap.}
Trap: Treating a link as healthy because it is up.\\
Better: Check corrected errors, bursts, retrains, and recovery time.

\paragraph{60-second answer.}
``Link-up is a state; health is a distribution over time. I read FEC activity for
consumed margin, look for bursts, and check recovery duration and repeats. I use
telemetry that separates hypotheses, compare affected and unaffected groups, and
judge severity from workload impact. Component FIT alone is not availability.''

\paragraph{Pressure.}
Pressure: ``Post-FEC looks clean. Why worry?''\\
Pivot: ``Clean post-FEC can still hide rising corrected-error demand and short
bursts. I would inspect those distributions before calling the link healthy.''

\footnotesize Depth: \Cref{ch:networking}.

\clearpage
\section{Chapter 11: Failure analysis}
\label{sec:cheat-ch11}

\textit{Goal.} Find the real cause without destroying the evidence, then stop it
from happening again.

\paragraph{Mental checklist.}
Protect $\rightarrow$ Preserve $\rightarrow$ Scope $\rightarrow$ Classify
$\rightarrow$ Hypothesize $\rightarrow$ Test $\rightarrow$ Prove $\rightarrow$
Fix $\rightarrow$ Prevent $\rightarrow$ Verify

\paragraph{What I need to know.}
\begin{itemize}
\item Do not reset, clean, or swap before preserving evidence.
\item A symptom is not a cause.
\item Timing helps: sudden, gradual, or intermittent.
\item Shared failures suggest a shared dependency.
\item Use reversible tests first.
\item A swap can isolate ownership without proving the physical mechanism.
\item It is acceptable to say the mechanism is not yet confirmed.
\item The case is not closed until recurrence control works on fresh data.
\end{itemize}

\paragraph{Key distinctions.}
\begin{center}
\footnotesize
\begin{tabular}{@{}ll@{}}
\toprule
Do not confuse & With \\
\midrule
Failure mode / symptom & Failure mechanism \\
Correction & Recurrence control \\
Ownership from a swap & Microscopic proof \\
\bottomrule
\end{tabular}
\end{center}

\paragraph{Context.}
\begin{center}
\footnotesize
\begin{tabular}{@{}ll@{}}
\toprule
Context & First priority \\
\midrule
Lab & Learn quickly and reproduce \\
Vetting or pilot lot & Protect lot integrity; collect representative evidence \\
Production fleet & Protect the workload; avoid disruptive tests until safe \\
\bottomrule
\end{tabular}
\end{center}
Same flow, different first priority.

\paragraph{Trap.}
Trap: Reseating first, or closing when the unit recovers.\\
Better: Preserve state, then close only after a recurrence control works on fresh
data.

\paragraph{60-second answer.}
``I start by protecting the workload and preserving the failing state. Then I
scope the issue by lane, module, host, lot, firmware, and site. I build a short
list of likely causes and use the smallest safe test, often a controlled swap, to
see what the problem follows. I do not call root cause from one correlation. Once
the cause is supported, I fix the immediate issue, put in a control to prevent
recurrence, and verify it on fresh data.''

\paragraph{Pressure.}
Pressure: ``What if it is in production and you cannot swap anything?''\\
Pivot: ``I protect the workload and collect non-disruptive evidence first. I
schedule intrusive tests only after traffic is drained or risk is contained.''

\footnotesize Depth: \Cref{ch:failure-modes,app:failure-analysis-reference}.

\clearpage
\section{Cross-chapter interview habits}
\label{sec:cheat-habits}

\paragraph{How to answer technical questions.}
\begin{enumerate}
\item Start with the decision or distinction.
\item State what you would check first.
\item Explain why that check is high value.
\item Name the main tradeoff.
\item Say what result would change your direction.
\item Avoid claiming more than the evidence supports.
\item End with the action or release decision.
\end{enumerate}

\paragraph{Useful phrases.}
\begin{itemize}
\item ``The first thing I would separate is\ldots''
\item ``I would not interpret that result until\ldots''
\item ``That raises a hypothesis, but it does not prove the cause.''
\item ``I would choose the cheapest safe test that separates the leading
      causes.''
\item ``The same measurement can support several claims, but the decision is
      different.''
\item ``I would protect the workload first, then preserve enough evidence to
      investigate.''
\item ``I would not call the case closed until the fix works on fresh data.''
\end{itemize}

\paragraph{Universal traps.}
\begin{itemize}
\item Jumping from correlation to cause.
\item Quoting a number without plane or condition.
\item Treating pass/fail as margin.
\item Treating a standard as the complete engineering argument.
\item Assuming a passing sample represents the whole population.
\item Treating link-up as health.
\item Treating final yield as process health.
\item Fixing the failed unit without preventing recurrence.
\end{itemize}

\textit{Design rule.} One page, one mental model, one spoken answer.

\clearpage
\small
\section{Chapters 7--11: 100 rapid-fire interview flashcards}
\label{sec:cheat-flashcards}

Use these as short memory prompts. Say the answer in one breath, then expand
only if asked. Depth:
\Cref{ch:product-readiness,ch:reliability,ch:manufacturing,ch:networking,ch:failure-modes}.

\subsection{Chapter 7: Product readiness}
\begin{enumerate}
\item \textbf{Q:} What is product readiness?\\
\textbf{A:} The full path from requirements to fleet learning.
\item \textbf{Q:} What is the first step?\\
\textbf{A:} Define measurable requirements and the release decision.
\item \textbf{Q:} Why review architecture early?\\
\textbf{A:} Catch thin margins before hardware makes changes expensive.
\item \textbf{Q:} What does bring-up prove?\\
\textbf{A:} The hardware, firmware, states, and basic link work reproducibly.
\item \textbf{Q:} What does characterization do?\\
\textbf{A:} Maps nominal behavior, spread, sensitivities, and failure edges.
\item \textbf{Q:} What does verification do?\\
\textbf{A:} Checks frozen requirements at named planes and conditions.
\item \textbf{Q:} What does system validation do?\\
\textbf{A:} Proves the product works in the intended system and workload.
\item \textbf{Q:} What does reliability qualification do?\\
\textbf{A:} Supports life and environmental claims.
\item \textbf{Q:} What does manufacturing validation do?\\
\textbf{A:} Proves the factory can build and control the design repeatedly.
\item \textbf{Q:} What is a pilot?\\
\textbf{A:} A bounded field experiment with telemetry, exit rules, and rollback.
\item \textbf{Q:} What justifies pilot expansion?\\
\textbf{A:} Metrics match the release model and no unexplained cohort appears.
\item \textbf{Q:} What makes a later test insufficient?\\
\textbf{A:} It cannot replace missing earlier evidence.
\item \textbf{Q:} What do EVT, DVT, and PVT mean?\\
\textbf{A:} Program phases, not universal proof categories.
\item \textbf{Q:} How do you interpret a phase gate?\\
\textbf{A:} Ask which build, which evidence, which risks, and which decision.
\item \textbf{Q:} Why not validate only on a reference host?\\
\textbf{A:} Production hosts can move the real failure edge.
\item \textbf{Q:} What must accompany every measurement?\\
\textbf{A:} Plane, metric, conditions, population, uncertainty, and decision.
\item \textbf{Q:} How do you choose the next test?\\
\textbf{A:} Pick the cheapest safe test that separates the leading causes.
\item \textbf{Q:} When is destructive analysis justified?\\
\textbf{A:} When non-destructive evidence cannot unblock a major decision.
\item \textbf{Q:} What do you protect when schedule is cut?\\
\textbf{A:} Requirements, measurement integrity, thin margins, and reversible
pilot.
\item \textbf{Q:} Best one-line readiness summary?\\
\textbf{A:} Define, prove, qualify, control, pilot, ramp, monitor, learn.
\end{enumerate}

\subsection{Chapter 8: Reliability qualification}
\begin{enumerate}
\setcounter{enumi}{20}
\item \textbf{Q:} What is reliability qualification?\\
\textbf{A:} A bounded confidence argument for life and environmental exposure.
\item \textbf{Q:} What is the first step in a qual plan?\\
\textbf{A:} Start with the claim, not the test list.
\item \textbf{Q:} What comes after the claim?\\
\textbf{A:} Identify the mechanism that could break it.
\item \textbf{Q:} What makes an accelerated stress useful?\\
\textbf{A:} It speeds up the same mechanism expected in the field.
\item \textbf{Q:} What makes a stress invalid?\\
\textbf{A:} It creates different failure physics.
\item \textbf{Q:} Why define observables before stress?\\
\textbf{A:} So acceptance is not invented after seeing the data.
\item \textbf{Q:} Why use representative samples?\\
\textbf{A:} One convenient lot does not represent the released population.
\item \textbf{Q:} What does zero failures prove?\\
\textbf{A:} A confidence bound, not a zero failure rate.
\item \textbf{Q:} What is FIT?\\
\textbf{A:} One failure per billion device-hours.
\item \textbf{Q:} What must accompany a FIT claim?\\
\textbf{A:} Population, exposure, failure definition, model, and confidence.
\item \textbf{Q:} What is DPPM?\\
\textbf{A:} Defective parts per million at a named quality boundary.
\item \textbf{Q:} FIT versus DPPM?\\
\textbf{A:} FIT uses time exposure; DPPM uses units inspected.
\item \textbf{Q:} HTOL versus burn-in?\\
\textbf{A:} HTOL supports life claims; burn-in is an optional production screen.
\item \textbf{Q:} Can burn-in replace qualification?\\
\textbf{A:} No.
\item \textbf{Q:} Can qualification replace screening?\\
\textbf{A:} No.
\item \textbf{Q:} What is a failure mode?\\
\textbf{A:} What you observe.
\item \textbf{Q:} What is a failure mechanism?\\
\textbf{A:} The physical or electrical cause.
\item \textbf{Q:} Why does failure timing matter?\\
\textbf{A:} Early, random, and late failures point to different risks.
\item \textbf{Q:} What is the role of standards like GR-468?\\
\textbf{A:} Common methods and language, not the whole engineering argument.
\item \textbf{Q:} Best one-line qual summary?\\
\textbf{A:} Claim, mechanism, stress, observable, sample, confidence, decision.
\end{enumerate}

\subsection{Chapter 9: Manufacturing validation}
\begin{enumerate}
\setcounter{enumi}{40}
\item \textbf{Q:} What does manufacturing validation prove?\\
\textbf{A:} The factory can build, measure, trace, and control the design.
\item \textbf{Q:} What do you freeze first?\\
\textbf{A:} Design, BOM, process, firmware, test software, and approved
suppliers.
\item \textbf{Q:} Why freeze the production reference?\\
\textbf{A:} So changes are visible and traceable.
\item \textbf{Q:} What is FAIR?\\
\textbf{A:} First Article Inspection Report.
\item \textbf{Q:} What does FAIR prove?\\
\textbf{A:} The first production-intent units match the released design and
process.
\item \textbf{Q:} What is MSA?\\
\textbf{A:} Measurement System Analysis.
\item \textbf{Q:} What is GR\&R?\\
\textbf{A:} Gauge repeatability and reproducibility.
\item \textbf{Q:} What does GR\&R ask?\\
\textbf{A:} Is the variation from the parts or from the measurement?
\item \textbf{Q:} What comes before yield interpretation?\\
\textbf{A:} Trust the measurement system.
\item \textbf{Q:} What is first-pass yield?\\
\textbf{A:} Units passing without retest or rework.
\item \textbf{Q:} Why is first-pass yield important?\\
\textbf{A:} It shows true process health.
\item \textbf{Q:} What is final yield?\\
\textbf{A:} Units eventually accepted after allowed rework and retest.
\item \textbf{Q:} Why can final yield mislead?\\
\textbf{A:} It can hide heavy rework and unstable process behavior.
\item \textbf{Q:} What is a Pareto chart?\\
\textbf{A:} A ranked view of the few causes driving most loss.
\item \textbf{Q:} Distribution versus capability?\\
\textbf{A:} Distribution is the raw shape; capability compares it with specs.
\item \textbf{Q:} What is process stability?\\
\textbf{A:} The process behaves predictably over time.
\item \textbf{Q:} Why must stability come before $C_{pk}$?\\
\textbf{A:} An unstable process makes capability numbers unreliable.
\item \textbf{Q:} $C_p$ versus $C_{pk}$?\\
\textbf{A:} $C_p$ is potential spread; $C_{pk}$ includes how off-center the
process is.
\item \textbf{Q:} Control limit versus spec limit?\\
\textbf{A:} Control limits describe process behavior; specs define product
acceptance.
\item \textbf{Q:} What is ATP?\\
\textbf{A:} The production acceptance test used to ship, hold, or reject units.
\item \textbf{Q:} How do you validate an ATP test?\\
\textbf{A:} Show it catches known bad cases, not just good units.
\item \textbf{Q:} What is a false accept?\\
\textbf{A:} A bad unit passes.
\item \textbf{Q:} What is a false reject?\\
\textbf{A:} A good unit fails.
\item \textbf{Q:} How do you set a production threshold?\\
\textbf{A:} Balance escapes, false rejects, cost, and customer risk.
\item \textbf{Q:} What is SPC?\\
\textbf{A:} Statistical Process Control for detecting drift over time.
\item \textbf{Q:} What must every reaction plan include?\\
\textbf{A:} Trigger, owner, containment, evidence, restart rule, and follow-up.
\item \textbf{Q:} Where should a defect be caught?\\
\textbf{A:} At the earliest reliable and economical point.
\item \textbf{Q:} How do you choose builds?\\
\textbf{A:} Cover real sources of variation and keep genealogy.
\item \textbf{Q:} Representation versus sample size?\\
\textbf{A:} First sample the right populations, then collect enough data.
\item \textbf{Q:} Best one-line manufacturing summary?\\
\textbf{A:} Freeze, trace, trust measurement, read yield, control, react, ramp.
\end{enumerate}

\subsection{Chapter 10: Optical links in operation}
\begin{enumerate}
\setcounter{enumi}{70}
\item \textbf{Q:} What is the chapter's main message?\\
\textbf{A:} Link-up is not the same as link health.
\item \textbf{Q:} What does pre-FEC error activity show?\\
\textbf{A:} The burden on the physical link before correction.
\item \textbf{Q:} What do corrected errors show?\\
\textbf{A:} FEC is working and margin may be shrinking.
\item \textbf{Q:} What do uncorrectable events show?\\
\textbf{A:} FEC could not recover the data.
\item \textbf{Q:} What is a retrain?\\
\textbf{A:} A recovery action that rebuilds link or lane state.
\item \textbf{Q:} Why can average BER mislead?\\
\textbf{A:} It can hide short, dangerous bursts.
\item \textbf{Q:} What should you ask about an error burst?\\
\textbf{A:} How big, how long, and what event it triggered.
\item \textbf{Q:} What is operational severity based on?\\
\textbf{A:} Workload consequence, not just optical symptoms.
\item \textbf{Q:} What does good recovery look like?\\
\textbf{A:} Fast, predictable, evidence-preserving, and not repetitive.
\item \textbf{Q:} Why is recovery part of product behavior?\\
\textbf{A:} A link will eventually hit faults and must return safely.
\item \textbf{Q:} What makes telemetry useful?\\
\textbf{A:} It separates causes, bounds a population, or triggers action.
\item \textbf{Q:} What is a cohort?\\
\textbf{A:} A group sharing lot, host, firmware, site, or another attribute.
\item \textbf{Q:} Why compare affected and unaffected cohorts?\\
\textbf{A:} To narrow the shared cause.
\item \textbf{Q:} Does correlation prove root cause?\\
\textbf{A:} No.
\item \textbf{Q:} Why can a rare event matter at fleet scale?\\
\textbf{A:} Large populations turn rare per-link events into frequent fleet
events.
\item \textbf{Q:} Why is FIT not availability?\\
\textbf{A:} Availability also depends on duration, redundancy, reroute, and
repair.
\item \textbf{Q:} What makes a pilot real?\\
\textbf{A:} Bounded exposure, telemetry, exit rules, and rollback.
\item \textbf{Q:} When do you pause deployment?\\
\textbf{A:} When risk is widening, unexplained, or hard to contain.
\item \textbf{Q:} What is the T1-to-T5 memory hook?\\
\textbf{A:} Trust measurement, know product, find edge, connect lab to system,
monitor fleet.
\item \textbf{Q:} Best one-line operations summary?\\
\textbf{A:} Watch margin, FEC, bursts, recovery, cohorts, and workload impact.
\end{enumerate}

\subsection{Chapter 11: Failure analysis}
\begin{enumerate}
\setcounter{enumi}{90}
\item \textbf{Q:} What is the first priority in production?\\
\textbf{A:} Protect the workload, then preserve evidence.
\item \textbf{Q:} What is the failure-analysis spine?\\
\textbf{A:} Protect, preserve, scope, classify, test, prove, fix, prevent,
verify.
\item \textbf{Q:} Why preserve before resetting?\\
\textbf{A:} Resets can erase the clues.
\item \textbf{Q:} What does scoping mean?\\
\textbf{A:} Find who is affected and what they share.
\item \textbf{Q:} What timing patterns matter?\\
\textbf{A:} Sudden, gradual, or intermittent.
\item \textbf{Q:} What is a hypothesis tree?\\
\textbf{A:} A short list of plausible causes and how to separate them.
\item \textbf{Q:} Why use reversible swaps first?\\
\textbf{A:} They localize ownership with low risk.
\item \textbf{Q:} If the fault follows the module, is root cause proven?\\
\textbf{A:} No. Ownership is localized; mechanism still needs proof.
\item \textbf{Q:} Correlation versus root cause?\\
\textbf{A:} Correlation moves together; root cause is reproducible and
stoppable.
\item \textbf{Q:} When is the case closed?\\
\textbf{A:} When the fix works on fresh data and recurrence is controlled.
\end{enumerate}

\subsection{Five universal fallback lines}
\begin{itemize}
\item ``I would start by clarifying the decision.''
\item ``I would not interpret that number without the plane and conditions.''
\item ``That raises a hypothesis, but it does not prove the cause.''
\item ``I would pick the smallest safe test that separates the leading causes.''
\item ``Based on the evidence, here is my call and next step.''
\end{itemize}

\normalsize
